% Source: Weinberg GR, physical PDF pp. 267--271
%         (printed pp. 244--248).
% Coverage: Section 9.9, Approximate Solutions to the Brans--Dicke Theory;
%           equations (9.9.1)--(9.9.29).
% Figures/tables/footnotes: no figures, tables, footnotes, or reference
%                           markers.
% Status: source-reviewed and compile-clean.
% Uncertainties: the Brans--Dicke scalar is renamed Phi_BD before
%                reparametrization; (9.9.3)--(9.9.4) are converted to the
%                target curvature sign; perturbative orders use
%                parenthesized superscripts.
%
% MODERNIZATION CHECK: The Brans--Dicke scalar is denoted Phi_BD before
% its reparametrization, avoiding a collision with the Newtonian potential
% phi. Covariant derivatives use nabla and Box. Perturbative-order numerals
% are rendered as parenthesized superscripts.
% MODERNIZATION CHECK: Equations (9.9.3) and (9.9.4) use the target
% curvature convention, so both sides are the negatives of the source
% equations. The component field equations (9.9.7)--(9.9.10) and all
% resulting metric solutions are consequently unchanged.

\subsection{Approximate Solutions to the Brans-Dicke Theory}\label{sec:9.9}

In order to test general relativity, it is useful to have in mind some
other theory with which to compare it. The Brans--Dicke theory described in
Section~\ref{sec:7.3} is identical with general relativity in the physical
interpretation of the metric \(g_{\mu\nu}\), and differs only in that a new
scalar field enters the gravitational field equations. In order to avoid
confusion with the Newtonian potential, we shall write the Brans--Dicke
scalar field as
\(\Phi_{\mathrm{BD}}=\mathcal G^{-1}(1+\xi)\), with \(\mathcal G\) a
constant of order \(G\), and \(\xi\) a scalar field defined by
\begin{equation}
  \Box\xi
  \equiv
  \nabla_\mu\nabla^\mu\xi
  =
  \frac{8\pi\mathcal G}{3+2\omega}T^\mu{}_\mu,
  \tag{9.9.1}\label{eq:9.9.1}
\end{equation}
\begin{equation}
  \xi\longrightarrow0
  \qquad\text{for}\qquad r\longrightarrow\infty.
  \tag{9.9.2}\label{eq:9.9.2}
\end{equation}
[See Eq.~\eqref{eq:7.3.13}.] We have dropped the subscript \(M\), but
\(T^{\mu\nu}\) should be understood as the energy-momentum tensor of
matter, excluding \(\xi\). Also, \(\omega\) is a dimensionless constant,
perhaps of order 6. The gravitational field equations are given by
Eq.~\eqref{eq:7.3.14}, in the target curvature convention, as
\begin{equation}
  \begin{split}
  R_{\mu\nu}-\frac12g_{\mu\nu}R
  ={}&
  8\pi\mathcal G(1+\xi)^{-1}T_{\mu\nu}\\
  &+\omega(1+\xi)^{-2}
  \left[
    (\nabla_\mu\xi)(\nabla_\nu\xi)
    -\frac12g_{\mu\nu}(\nabla_\rho\xi)(\nabla^\rho\xi)
  \right]\\
  &+(1+\xi)^{-1}
  \left[
    \nabla_\mu\nabla_\nu\xi-g_{\mu\nu}\Box\xi
  \right].
  \end{split}
  \tag{9.9.3}\label{eq:9.9.3}
\end{equation}
By using Eq.~\eqref{eq:9.9.1} to determine \(\Box\xi\), and contracting
Eq.~\eqref{eq:9.9.3} to find \(R\), we can rewrite this in the form
\begin{equation}
  \begin{split}
  R_{\mu\nu}={}&
  8\pi\mathcal G(1+\xi)^{-1}
  \left[
    T_{\mu\nu}
    -g_{\mu\nu}T^\lambda{}_\lambda
     \left(\frac{\omega+1}{2\omega+3}\right)
  \right]\\
  &+\omega(1+\xi)^{-2}
    (\nabla_\mu\xi)(\nabla_\nu\xi)
  +(1+\xi)^{-1}\nabla_\mu\nabla_\nu\xi.
  \end{split}
  \tag{9.9.4}\label{eq:9.9.4}
\end{equation}

It follows from Eqs.~\eqref{eq:9.9.1} and \eqref{eq:9.9.2} that \(\xi\)
has the expansion
\begin{equation}
  \xi=\xi^{(2)}+\xi^{(4)}+\cdots,
  \tag{9.9.5}\label{eq:9.9.5}
\end{equation}
where \(\xi^{(N)}\) is of order \(\bar v^N\), and in particular
\begin{equation}
  \nabla^2\xi^{(2)}
  =
  -\frac{8\pi\mathcal G}{3+2\omega}T^{00(0)}.
  \tag{9.9.6}\label{eq:9.9.6}
\end{equation}
The field equations are now given by
Eqs.~\eqref{eq:9.9.4}--\eqref{eq:9.9.6} and
\eqref{eq:9.1.37}--\eqref{eq:9.1.40} as
\begin{equation}
  \nabla^2g_{00}^{(2)}
  =
  -8\pi\mathcal G
  \left(\frac{2\omega+4}{2\omega+3}\right)T^{00(0)},
  \tag{9.9.7}\label{eq:9.9.7}
\end{equation}
\begin{equation}
  \begin{split}
  \nabla^2g_{00}^{(4)}
  ={}&
  \partial_t^2g_{00}^{(2)}
  +g_{ij}^{(2)}\partial_i\partial_jg_{00}^{(2)}
  -\bigl(\boldsymbol\nabla g_{00}^{(2)}\bigr)^2\\
  &+8\pi\mathcal G
   \left(\frac{2\omega+4}{2\omega+3}\right)^2
   \xi^{(2)}T^{00(0)}\\
  &-8\pi\mathcal G
   \left(\frac{2\omega+2}{2\omega+3}\right)T^{ii(2)}\\
  &+16\pi\mathcal G g_{00}^{(2)}T^{00(0)}
   \left(\frac{2\omega+4}{2\omega+3}\right)\\
  &-8\pi\mathcal G
   \left(\frac{2\omega+4}{2\omega+3}\right)T^{00(2)}
  -2\omega\bigl(\partial_t\xi^{(2)}\bigr)^2\\
  &-2\partial_t^2\xi^{(2)}
  +2\Gamma^{i(2)}{}_{00}\partial_i\xi^{(2)}.
  \end{split}
  \tag{9.9.8}\label{eq:9.9.8}
\end{equation}
\begin{equation}
  \nabla^2g_{i0}^{(3)}
  =
  16\pi\mathcal G T^{i0(1)}
  -2\partial_i\partial_t\xi^{(2)},
  \tag{9.9.9}\label{eq:9.9.9}
\end{equation}
\begin{equation}
  \nabla^2g_{ij}^{(2)}
  =
  -8\pi\mathcal G T^{00(0)}\delta_{ij}
   \left(\frac{2\omega+2}{2\omega+3}\right)
  -2\partial_i\partial_j\xi^{(2)}.
  \tag{9.9.10}\label{eq:9.9.10}
\end{equation}

From Eq.~\eqref{eq:9.9.7} it follows that the gravitational constant
measured by observation of slowly moving particles or in time-dilation
experiments is not \(\mathcal G\), but rather is
\begin{equation}
  G=
  \left(\frac{2\omega+4}{2\omega+3}\right)\mathcal G.
  \tag{9.9.11}\label{eq:9.9.11}
\end{equation}
That is, we have the usual relation between \(g_{00}^{(2)}\) and the
Newtonian potential \(\phi\),
\begin{equation}
  g_{00}^{(2)}=-2\phi,
  \tag{9.9.12}\label{eq:9.9.12}
\end{equation}
provided we define \(\phi\) by
\begin{equation}
  \nabla^2\phi=4\pi G T^{00(0)}.
  \tag{9.9.13}\label{eq:9.9.13}
\end{equation}
Also, it follows from Eqs.~\eqref{eq:9.9.6} and \eqref{eq:9.9.13} that
\[
  \xi^{(2)}=-(\omega+2)^{-1}\phi,
\]
and the field equations for \(g_{00}^{(4)},g_{i0}^{(3)}\), and
\(g_{ij}^{(2)}\) are
\begin{equation}
  \begin{split}
  \nabla^2g_{00}^{(4)}
  ={}&
  -2\left(\frac{\omega+1}{\omega+2}\right)\partial_t^2\phi
  -2g_{ij}^{(2)}\partial_i\partial_j\phi
  -2\left(\frac{2\omega+5}{\omega+2}\right)
   (\boldsymbol\nabla\phi)^2\\
  &-8\pi G\left(4+\frac1{\omega+2}\right)
   \phi T^{00(0)}
  -8\pi G\left(\frac{2\omega+2}{2\omega+4}\right)^2T^{ii(2)}\\
  &-8\pi G T^{00(2)}
  -\frac{2\omega}{(\omega+2)^2}(\partial_t\phi)^2.
  \end{split}
  \tag{9.9.14}\label{eq:9.9.14}
\end{equation}
\begin{equation}
  \nabla^2g_{i0}^{(3)}
  =
  16\pi G
  \left(\frac{2\omega+3}{2\omega+4}\right)T^{i0(1)}
  +\frac{2}{\omega+2}\partial_i\partial_t\phi,
  \tag{9.9.15}\label{eq:9.9.15}
\end{equation}
\begin{equation}
  \nabla^2g_{ij}^{(2)}
  =
  -8\pi G T^{00(0)}\delta_{ij}
   \left(\frac{\omega+1}{\omega+2}\right)
  +\frac{2}{\omega+2}\partial_i\partial_j\phi.
  \tag{9.9.16}\label{eq:9.9.16}
\end{equation}

As an example, let us consider the field of a static spherically symmetric
mass. The Newtonian potential is then a function of \(r\) alone, and
Eq.~\eqref{eq:9.9.16} gives
\begin{equation}
  \begin{split}
  g_{ij}^{(2)}
  ={}&
  -2\delta_{ij}\left(\frac{\omega+1}{\omega+2}\right)\phi\\
  &+\frac{2}{\omega+2}
  \left\{
    \left(\delta_{ij}-\frac{3x_ix_j}{r^2}\right)
    \frac1{r^3}\int_0^r r'^2\phi(r')\,\dd r'
    +\frac{x_ix_j}{r^2}\phi
  \right\}.
  \end{split}
  \tag{9.9.17}\label{eq:9.9.17}
\end{equation}
Outside the mass, we have
\begin{equation}
  \phi=-\frac{MG}{r},
  \tag{9.9.18}\label{eq:9.9.18}
\end{equation}
and so Eq.~\eqref{eq:9.9.17} gives
\begin{equation}
  \begin{split}
  g_{ij}^{(2)}
  ={}&
  \left(\frac{2\omega+1}{\omega+2}\right)
  \frac{MG}{r}\delta_{ij}
  +\frac{MG}{\omega+2}\frac{x_ix_j}{r^3}\\
  &+\frac{2MGR^2}{\omega+2}
  \left(\delta_{ij}-\frac{3x_ix_j}{r^2}\right)\frac1{r^3},
  \end{split}
  \tag{9.9.19}\label{eq:9.9.19}
\end{equation}
where \(R\) is an effective radius, defined by
\begin{equation}
  MGR^2
  \equiv
  \int_0^\infty
  \left[\phi(r)+\frac{MG}{r}\right]r^2\,\dd r.
  \tag{9.9.20}\label{eq:9.9.20}
\end{equation}
(The integrand vanishes outside the mass, so we are free to change its
upper limit from \(r\) to \(\infty\).) Using Eqs.~\eqref{eq:9.9.18} and
\eqref{eq:9.9.19} in Eq.~\eqref{eq:9.9.14} gives
\[
  \nabla^2g_{00}^{(4)}
  =
  -\frac{2(2\omega+3)M^2G^2}{(\omega+2)r^4}
  -\frac{24M^2G^2R^2}{(\omega+2)r^6}.
\]
The solution is
\begin{equation}
  g_{00}^{(4)}
  =
  -\frac{(2\omega+3)M^2G^2}{(\omega+2)r^2}
  -\frac{2M^2G^2R^2}{(\omega+2)r^4}
  +\frac{\kappa M^2G^2}{rR},
  \tag{9.9.21}\label{eq:9.9.21}
\end{equation}
where \(\kappa\) is a dimensionless constant, which must be determined by
the condition that the exterior solution \eqref{eq:9.9.21} fit smoothly
onto a nonsingular internal solution.

The results \eqref{eq:9.9.19}--\eqref{eq:9.9.21} make it appear that the
gravitational field outside a spherical static mass depends upon the size
and distribution of the mass. However, this size-dependent effect can be
eliminated by a suitable redefinition of \(M\) and \(\mathbf x\):
\begin{equation}
  M'
  =
  M\left(1-\frac{\kappa MG}{R}\right),
  \tag{9.9.22}\label{eq:9.9.22}
\end{equation}
\begin{equation}
  \mathbf x'
  =
  \mathbf x\left[
    1+\frac{MGR^2}{(\omega+2)r^3}
  \right].
  \tag{9.9.23}\label{eq:9.9.23}
\end{equation}
The last two terms in Eq.~\eqref{eq:9.9.21} and the last term in
Eq.~\eqref{eq:9.9.19} are then canceled by the changes in
\(g_{00}^{(2)}\) and \(g_{ij}^{(0)}\), and so, dropping primes, we now have
\begin{equation}
  g_{00}^{(2)}=\frac{2MG}{r},
  \tag{9.9.24}\label{eq:9.9.24}
\end{equation}
\begin{equation}
  g_{00}^{(4)}
  =
  -\frac{(2\omega+3)M^2G^2}{(\omega+2)r^2},
  \tag{9.9.25}\label{eq:9.9.25}
\end{equation}
\begin{equation}
  g_{ij}^{(2)}
  =
  \left(\frac{2\omega+1}{\omega+2}\right)
  \frac{MG}{r}\delta_{ij}
  +\frac{MG}{\omega+2}\frac{x_ix_j}{r^3}.
  \tag{9.9.26}\label{eq:9.9.26}
\end{equation}
Thus the Brans--Dicke theory shares the property of the Einstein theory
that the gravitational field outside a static spherically symmetric mass
depends on \(M\), but not on any other property of the mass.

This solution may be compared with the general Robertson expansion
\eqref{eq:8.3.7} in harmonic coordinates, which (with \(\alpha=1\)) gives
\[
  g_{00}^{(2)}=\frac{2MG}{r},
\]
\[
  g_{00}^{(4)}
  =
  -(\gamma-1+2\beta)\frac{M^2G^2}{r^2},
\]
\[
  g_{ij}^{(2)}
  =
  (3\gamma-1)\delta_{ij}\frac{MG}{r}
  +(1-\gamma)\frac{MGx_ix_j}{r^3}.
\]
Thus the Brans--Dicke results
\eqref{eq:9.9.24}--\eqref{eq:9.9.26} can be summarized by giving formulas
for the Robertson parameters:
\begin{equation}
  \gamma=\frac{\omega+1}{\omega+2},
  \qquad
  \beta=1.
  \tag{9.9.27}\label{eq:9.9.27}
\end{equation}
These formulas were already used in the last chapter to compare the
Brans--Dicke theory with experiment.

We also note that the element \(g_{i0}^{(3)}\equiv\zeta_i\) of the metric
tensor is given for a static system by Eq.~\eqref{eq:9.9.15} as
\begin{equation}
  \zeta_i(\mathbf x,t)
  =
  -4G\left(\frac{2\omega+3}{2\omega+4}\right)
  \int
  \frac{T^{i0(1)}(\mathbf x',t)}
       {|\mathbf x-\mathbf x'|}\,\dd^3x'.
  \tag{9.9.28}\label{eq:9.9.28}
\end{equation}
Thus the effects of the rotation of a spherical mass on the precession of
spins and perihelia are smaller in the Brans--Dicke theory (for
\(0<\omega<\infty\)) than in general relativity, by a factor
\((2\omega+3)/(2\omega+4)\).

By far the most dramatic tests of the Brans--Dicke theory are those that
also test the ``very strong'' Principle of Equivalence. At any point \(P\)
in a gravitational field we can choose a locally inertial coordinate
system, for which \(g_{\mu\nu}=\eta_{\mu\nu}\) and
\(\Gamma^\lambda{}_{\mu\nu}=0\) at that point. However, the Brans--Dicke
field \(\xi\) is a scalar, and hence will not vanish at \(P\), being given
by Eqs.~\eqref{eq:9.9.6} and \eqref{eq:9.9.13} as
\[
  \xi\simeq\xi^{(2)}=-(\omega+2)^{-1}\phi,
\]
where \(\phi\) is the Newtonian gravitational potential.
Equation~\eqref{eq:9.9.4} shows that in this coordinate system the
gravitational field of a small mass introduced at \(P\) can be calculated
as usual, but with the gravitational constant \(G\) replaced with
\begin{equation}
  G_{\mathrm{eff}}
  =
  G(1+\xi)^{-1}
  \simeq
  G\left[1+(\omega+2)^{-1}\phi\right].
  \tag{9.9.29}\label{eq:9.9.29}
\end{equation}
For instance, with \(\omega=6\) and \(\phi\) at the surface of the earth
equal to \(-6.9\times10^{-10}\), the effective gravitational constant
measured by Cavendish experiments on the surface of the earth is
\emph{smaller} than the ``true'' gravitational coupling constant that
would be measured on a satellite in a high orbit by a factor
\(1-8\times10^{-11}\).
